\uuid{9Y0R}
\exo7id{4749}
\titre{exo7 4749}
\auteur{quercia}
\organisation{exo7}
\datecreate{2010-03-16}
\isIndication{false}
\isCorrection{false}
\chapitre{Topologie}
\sousChapitre{Topologie des espaces vectoriels normés}
\cpge{mpsi-mp2i.calcul-matriciel-systemes.e, pcsi.determinants.c, ptsi.determinants.a, bcpst1.matrices, mp-mpi.topologie-espaces-normes.b, pc.espaces-vectoriels-normes.b, psi.espaces-vectoriels-normes.b}
\module{Analyse}
\niveau{L2}
\difficulte{}

\contenu{
\question{
Soit $B$ une matrice antisym{\'e}trique.
On suppose que la suite $(B^n)$ converge vers une matrice~$C$.
Que peut-on dire de $C$~?
}
}
